\section{Innledning}
I denne rapporten vil halvlederdiodens virkemåte i vekselspenningskretser bli undersøkt gjennom både teoretiske beregninger og praktiske målinger. Arbeidet vil ta for seg hvordan vanlige silisiumdioder og zenerdioder oppfører seg i klippekretser, enveis likeretter, toveis likeretter og toveis likeretter med RC-ledd. Hensikten vil være å skape forståelse for hvordan dioder kan brukes til å forme signaler, likerette vekselspenning og påvirke nivået på DC-spenning og ripple i en krets.

Rapporten vil særlig søke å besvare hvordan spenningskurver endres når dioder kobles i ulike konfigurasjoner, hvordan DC-spenning kan beregnes ved enveis og toveis likeretting, og hvordan en kondensator i et RC-ledd vil påvirke utgangsspenningen. Det vil også bli undersøkt i hvilken grad de praktiske målingene stemmer overens med de teoretiske beregningene, og hvilke årsaker som kan forklare eventuelle avvik.

Rapportens omfang vil være avgrenset til kretsene og målingene som inngår i selve oppgaven. Det vil bli benyttet tilgjengelig laboratorieutstyr i form av oscilloskop, digitalt multimeter, signalgenerator og transformator, samt de komponentene som er oppgitt i komponentlisten. Videre vil analysene i hovedsak bygge på praktisk diode-modell, der silisiumdiodens spenningsfall tilnærmes til $0.7V$. Målingene vil også bli gjennomført med de gitte betingelsene for signalgenerator og oppkobling slik oppgaven beskriver. 

Rapporten vil ikke gi en fullstendig analyse av alle ikke-ideelle forhold i halvlederdioder, men vil være begrenset til det som kan observeres og måles med tilgjengelige instrumenter. Små avvik vil derfor kunne oppstå som følge av komponenttoleranser, måleusikkerhet, variasjoner i spenningsfall over dioder og andre forhold i oppkoblingen av kretsene.

På bakgrunn av dette vil rapporten først presentere nødvendig teori, deretter forhåndsberegninger, gjennomføring og måleresultater, før resultatene og avvik til slutt vil bli diskutert. Målet vil være å vise sammenhengen mellom teori og praksis, og å belyse hvordan dioder er sentralt i elektroniske komponenter og anvendelser knyttet til klipping og likeretting av signaler.
